\documentclass[11pt]{article}
\usepackage[margin=25mm]{geometry}
\usepackage{fontspec}\setmainfont{Latin Modern Roman}
\usepackage{amsmath,amssymb,mathtools}
\usepackage{unicode-math}\setmathfont{Latin Modern Math}
\usepackage{xurl,hyperref}
\hypersetup{colorlinks=true,urlcolor=blue}
\setcounter{secnumdepth}{0}
\setlength{\parindent}{0pt}\setlength{\parskip}{5pt}
\setlength{\emergencystretch}{3em}
\newcommand{\sref}[2]{\hyperlink{source:#1:#2}{[S#2]}}
\newcommand{\cataloguescope}[1]{\paragraph{Recorded target.}#1}
\begin{document}
\section{Kourovka Notebook Problem 21.148}\label{statement}

\subsection{Definitions and mathematical statement}

\paragraph{Shared notation.}$\mathbb N=\{1,2,\ldots\}$, and $\mathbb Z,\mathbb Q,\mathbb R,\mathbb C$ denote the integers, rationals, reals and complex numbers. Any other conventions are specified locally.


A right order on $G$ is a total order $<$ satisfying $x<y\Rightarrow xg<yg$ for all $g\in G$. A subgroup $H$ is convex for this order if $h_1<x<h_2$, with $h_1,h_2\in H$, implies $x\in H$. Let $\mathrm{RC}(G)$ consist of subgroups convex for at least one right order; different subgroups may use different orders. Order $\mathrm{RC}(G)$ by inclusion. Write $\wedge,\vee$ for the intrinsic meet and join in this lattice. Distributivity means both $H\wedge(K\vee L)=(H\wedge K)\vee(H\wedge L)$ and its dual; a chain means any two elements are comparable.
\paragraph{Statement recorded in the supplied source.}
For every right-orderable $G$, is the lattice $\mathrm{RC}(G)$ distributive if and only if it is a chain? \sref{statement}{1}\sref{statement}{2}
\subsection{Short English statement}
Is distributivity of the lattice of right-relatively convex subgroups equivalent to their being totally ordered by inclusion?
\subsection{Sources}
\begin{description}
\item[\hypertarget{source:statement:1}{S1}] ULAM.AI and the UnsolvedMath Contributors, \emph{UnsolvedMath: A Curated Collection of Open Mathematics Problems}, catalogue entry. CC BY 4.0. \url{https://huggingface.co/datasets/ulamai/UnsolvedMath}.
\item[\hypertarget{source:statement:2}{S2}] E. I. Khukhro and V. D. Mazurov (editors), The Kourovka Notebook, Unsolved Problems in Group Theory, No.21, version48 (6 October2026), Problem 21.148 and its remarks. \url{https://arxiv.org/abs/1401.0300v48}.
\end{description}
\label{end:statement}
\end{document}
